\section{Introduction}
\label{sec:intro}

Many economic decisions change which decisions come next. In Monopoly, selling the houses on a color group makes its deeds mortgageable. Selling the last deed of a group to a rival lets that rival build. Lending to a nearly bankrupt opponent keeps a blocker in play, and refusing to lend can deliver the same deed to whoever is owed money. An agent that plays well has to say, at least implicitly, what an action will make possible for itself and for others, and then act on it.

Language-model agents can be asked to write that reasoning down. In MonopolyBench each response carries a public message and, under the standard prompt, a private note attached to the applied action, and most providers also return a reasoning text. These records let us ask a question that outcomes alone cannot answer. When an agent states, before acting, how an action will change its own or another player's options, does what it then does agree with what it said? This is narrower than asking whether the agent reasons well, and the records can answer it.

We study twenty complete four-player games played by GPT-6.1 Sol, Gemini 4 Argon, Claude Opus 5.5 and Grok 4.7 through one gateway. The corpus contains 6,989 decisions, 920 responses in 526 negotiation threads, and 188 debts that required a liquidation menu. Every offer is linked to the state its recipient saw and to its response, and every debt to its first menu, applied actions and outcome.

We grade each result against a five-step evidence ladder. A statement can mention a consequence for some player (\Level{stated}). The consequence can match the state and rules (\Level{accurate}). The applied action can agree with it (\Level{consistent}). The action can change when the relevant opportunity changes (\Level{sensitive}). And the behavior can improve an outcome against a reference (\Level{effective}). Observational records can reach the third step. They cannot reach the fourth or fifth without a comparison that holds other things fixed, and this study runs no new games, replays or model queries.

\paragraph{Own options.} The clearest evidence concerns the agent's own changing menu. In sixteen debts, cash, the salvage value of every building and every mortgage listed on the first menu together fall short. Payment then requires mortgages that appear only after buildings are sold. Agents pay fourteen of these debts and declare bankruptcy in two. Before the first action, the recorded text names the released mortgages, or a budget that includes them, in thirteen of the fourteen payments and raises them as an open question in the fourteenth. In both bankruptcies the text adds up only the listed items and the agent declares at once. The worked example is a ten-action liquidation whose announced remainder of \$14 is exactly what remains.

\paragraph{Rivals' options.} Bargaining evidence is broader and weaker. Recipients often refuse or reprice offers that would complete another player's color group, and their notes usually cite what the proposer could build next. In opening offers, swaps that complete only the proposer's group are rejected more often than swaps that complete no group (95/116 versus 37/72). Cash bids show no corresponding increase (42/70 versus 67/103). The swap gap keeps its sign in every proposer stratum with data, but its size depends on who proposed the comparison offers. Without Grok-authored offers it falls from 30.5 to 16.7 percentage points, and a resampling interval that respects the seed design then includes zero. Recipients' stated reasons cite completion and building under both offer forms, so the reasons do not explain why only swaps differ. We report this pattern as exploratory and do not interpret it as an identified response to the rival's next move.

\paragraph{Conduct in context.} Eight incidents show these statements in play: refusing and pricing against a rival's building cash, a factual correction, a settled bluff, a selective claim about financing, a public appeal, a withdrawn assurance and the ten-action payment. Each comes with exact quotations, named channels, checked delivery where claimed, and its recorded continuation.

The paper makes three contributions.
\begin{enumerate}
\item A complete debt census that separates easy from dependent liquidation and pairs each dependent case with the prospective text that preceded it (Section~\ref{sec:finance}).
\item A complete opening-offer census by form, completion, proposer, recipient and seed block, with the composition and uncertainty a reader needs to judge the form pattern (Section~\ref{sec:bargaining}).
\item Eight contextual incidents with verified quotations, channels and continuations, each assigned the evidence level it reaches (Sections~\ref{sec:finance}--\ref{sec:conduct}).
\end{enumerate}

The study describes four configured systems under unequal conditions. Opus played most decisions through a refusal-recovery path without private notes, and Grok ran with forced tool choice and sometimes returned text naming a different action from the call the engine applied. The results are observations of this cohort, not rankings of models. Appendix~\ref{app:questions} answers the eight motivating research questions and three broader ones at the level the records support.
